\chapter{Introducción}
\label{chap:introduccion}

El problema de la generación automática de horarios académicos es un desafío clásico de
optimización combinatoria en instituciones de educación superior. En la Facultad Politécnica de la
Universidad Nacional de Asunción, la asignación eficiente de recursos requiere equilibrar múltiples
restricciones operativas y pedagógicas.
~\cite{prime-number-theorem}

\section{Motivación y planteamiento del problema}

Lorem ipsum dolor sit amet, consectetur adipiscing elit. Sed do eiusmod tempor incididunt ut labore
et dolore magna aliqua. Ut enim ad minim veniam, quis nostrud exercitation ullamco laboris nisi ut
aliquip ex ea commodo consequat.
~\cite{miticagenda}

Duis aute irure dolor in reprehenderit in voluptate velit esse cillum dolore eu fugiat nulla
pariatur. Excepteur sint occaecat cupidatat non proident, sunt in culpa qui officia deserunt mollit
anim id est laborum:

\begin{itemize}
    \item **Restricciones duras:** Lorem ipsum dolor sit amet, consetetur sadipscing elitr, sed
        diam nonumy eirmod tempor invidunt.
    \item **Restricciones blandas:** At vero eos et accusam et justo duo dolores et ea rebum, stet
        clita kasd gubergren.
\end{itemize}

\section{Objetivos}

\subsection{Objetivo general}

Desarrollar y evaluar un marco algorítmico comparativo para la generación automática de horarios
universitarios adaptado a la estructura curricular de la FP-UNA.

\subsection{Objetivos específicos}

\begin{enumerate}
    \item Lorem ipsum dolor sit amet, consetetur sadipscing elitr, sed diam nonumy eirmod tempor
        invidunt ut labore et dolore.
    \item Takimata sanctus est lorem ipsum dolor sit amet. Lorem ipsum dolor sit amet, consetetur
        sadipscing elitr.
    \item Duis autem vel eum iriure dolor in hendrerit in vulputate velit esse molestie consequat.
\end{enumerate}

\section{Estructura del documento}

El presente documento de tesis se organiza de la siguiente manera:

En el **Capítulo 1**, lorem ipsum dolor sit amet, consectetur adipiscing elit, sed do eiusmod
tempor incididunt ut labore. En el **Capítulo 2**, se presentan los detalles del diseño
experimental. Finalmente, los datos completos y las métricas detalladas se incluyen en el **Apéndice A**.
